% =========================================================================
%  Resumo
% =========================================================================

\chapter*{Resumo}
\addcontentsline{toc}{chapter}{Resumo}

\noindent
% --- INSTRUÇÕES (apagar antes da entrega) ---------------------
% O resumo deve ter entre 150 e 250 palavras e responder a:
%   1. Qual o problema que o projeto resolve?
%   2. Qual a abordagem técnica escolhida?
%   3. Quais os principais resultados ou contribuições?
% Evitar entrar em detalhes técnicos extensos — esses ficam para
% os capítulos seguintes.
% ---------------------------------------------------------------

[Substituir pelo texto do resumo. Apresentar em poucas linhas o
contexto do projeto GREENHERB, os principais requisitos abordados, a
arquitetura geral adotada (browser, API, base de dados) e os resultados
relevantes. O leitor que apenas leia este resumo deve ficar com uma
ideia clara do que foi feito e do que se demonstrou.]

\vspace{0.5cm}
\noindent\textbf{Palavras-chave:} aplicação web, gestão de estufa, Node.js,
MongoDB, armazenamento no browser, JWT, OpenAPI.

\clearpage
